\section{FunQA Dataset}

% \begin{figure}[htbp]
%     \centering
%     \includegraphics[scale=0.5]{figures/dataset_main.pdf}
%     \caption{Pipleline of FunQA Dataset.}
%     \label{fig:dataset_main}
% \end{figure}

In this section, we provide a detailed explanation of the design principles that guided the creation of the FunQA dataset and its subsets. Additionally, we introduce our novel VideoQA tasks tailored for FunQA, and FunQA data statistics in Figure~\ref{fig:dataset statistics}. We introduce our construction pipeline in the end.


\begin{figure}[!t]

\begin{subfigure}[b]{0.3\textwidth}
    \centering
    \includegraphics[height=3cm]{figures/fig_word_cloud_h.png}
    \caption{Word cloud of HumorQA}
    \label{fig:wc_h}
  \end{subfigure}
  \hfill
  \begin{subfigure}[b]{0.3\textwidth}
    \centering
    \includegraphics[height=3cm]{figures/fig_word_cloud_c.png}
    \caption{Word cloud of CreativeQA}
    \label{fig:wc_c}
  \end{subfigure}
  \hfill
  \begin{subfigure}[b]{0.3\textwidth}
    \centering
    \includegraphics[height=3cm]{figures/fig_word_cloud_m.png}
    \caption{Word cloud of MagicQA}
    \label{fig:wc_m}
  \end{subfigure}

  \begin{subfigure}[b]{0.3\textwidth}
    \centering
    \includegraphics[height=3cm]{figures/fig_split_big.pdf}
    \caption{Number of videos and split}
    \label{fig:fun_split}
  \end{subfigure}
  \hfill  
 \begin{subfigure}[b]{0.3\textwidth}
    \centering
    \includegraphics[height=2.6cm]{figures/fig_video_len_big.pdf}
    \caption{Total video length}
    \label{fig:fun_len}
  \end{subfigure}
  \hfill
  \begin{subfigure}[b]{0.3\textwidth}
    \centering
    \includegraphics[height=3cm]{figures/fig_vqa_num_big.pdf}
    \caption{Number of QA pairs}
    \label{fig:fun_qa}
  \end{subfigure}
  \hfill


  
  \begin{subfigure}[b]{0.3\textwidth}
    \centering
    \includegraphics[height=3cm]{figures/fig_moment_distrib.png}
    \caption{Time Dist. of localization task}
    \label{fig:fun_moment}
  \end{subfigure}
  \hfill
  \begin{subfigure}[b]{0.35\textwidth}
    \centering
    \includegraphics[height=3cm]{figures/fig_avg_word_len.pdf}
    \caption{Avg. word count of answers}
    \label{fig:fun_avg}
  \end{subfigure}
  \hfill
  \begin{subfigure}[b]{0.3\textwidth}
    \centering
    \includegraphics[height=3cm]{figures/fig_val.pdf}
    \caption{Anno. consensus evaluation}
    \label{fig:fun_val}
  \end{subfigure}
  \hfill
  

  

  \caption{\textbf{Statistics of FunQA Dataset.} FunQA consists of three subsets, each corresponding to different video types, and is annotated with free-text QA pairs. The first row displays word clouds representing critical annotations for each subset. The second row provides key dataset statistics, including the number of videos for different splits, video length, and QA pair count for three subsets. In the last row, (g) highlights the high-frequency time span of the answer for localization questions in red, (h) shows the average word count of answers, and (i) presents the percentage of consensus between annotators for the same answer in a sampled set.
  }
  \vspace{-0.4cm}
  \label{fig:dataset statistics}
\end{figure}


\subsection{Video Selection}
% \noindent\textbf{Design principle}\quad 
In constructing the dataset, we adhered to three principles to address the challenges in video understanding capabilities: our dataset, FunQA, is \textbf{visual centered} and emphasizes \textbf{counter-intuitive reasoning}, \textbf{spatial-temporal reasoning}. Based on these principles, we collect 4365 videos from 3 different art genres and created three subsets: HumorQA, CreativeQA, and MagicQA.\\
% \jingkang{Based on the principle, we collect xxx videos from 3 different art genres....... What is your subset, why these datasets, simple stats of these dataset.}
\noindent\textbf{HumorQA}\quad
HumorQA composed of 1,769 meticulously curated web videos, serves as a unique source of insight into human humor comprehension. Notably, it contains the shortest average video length of 7s among the three subsets. We believe that the human process of understanding humor is complex and deep, requiring a holistic understanding of the video and adding a degree of common sense to it. Psychological research has demonstrated that humor arises from the incongruity \cite{kant1987critique,sep-humor} between reality and expectations, flourishing with the skillful juxtaposition and transformation of events \cite{latta1999basic,boyd2004laughter,koestler2020act}. This makes humorous videos a valuable asset for the VideoQA dataset, anticipated to enhance a model's proficiency in integrating information and performing deep reasoning.\\
\noindent\textbf{CreativeQA}\quad 
CreativeQA is a collection of 927 videos averaging 48s in length from a TV show called Kasou Taishou \cite{ntv}. This program, showcasing original and novel skits performed by various amateur groups and judged by a panel, boasts a strong creative flair 
 \cite{runco2012standard}. The essence of the show lies in using a mix of people and props to mimic reality, with audiences deriving pleasure from information integration and comparison. We anticipate that the imitation nature of the show will challenge the model's capacity for information extraction, while the longer video length and need for understanding creativity will put to test the model's comprehension of spatial-temporal information.\\
\noindent\textbf{MagicQA}\quad
MagicQA encapsulates 1672 magic performance videos sourced from across the web, spanning various genres like camera magic, close-up magic, and stage magic. The essence of magic revolves around the creation of seemingly impossible illusions \cite{nelms2012magic}, employing diverse effects such as disappearance, creation, and transformation. These illusions are infused with abundant spatial-temporal information. Through this dataset, we aim to empower the model to not only track the ensuing changes in objects but also unravel the underlying mechanics \cite{lamont2005magic} of these transformations.
% \noindent\textbf{FunQA}\quad
% In general, our meticulous filtering process ensures that the information within the aforementioned three types of art remains focused on visual content. These art forms naturally encompass complex spatial-temporal information, and their appreciation also requires deep human reasoning. Therefore, they align perfectly with our design guidelines.
\subsection{Task Definition}
% \jingkang{Each datasets has the following 4 types of tasks.....}


To comprehensively evaluate the model's ability to understand surprising videos, we designed the following 4 types of tasks for each subset: \\
\noindent\textbf{Counter-intuitive Timestamp Localization Task}\quad
The localization task is the base task to assess the model's comprehension abilities. It involves localizing counter-intuitive segments within the video, answers expressed in either seconds or frames. This task serves as the basis for the subsequent two main tasks in the three subsets, where the focus shifts to locating moments of humor, creativity, and magical effects, respectively. Successfully completing this task demands the model's understanding of the video's overall content, incorporating both temporal and spatial information.\\
\noindent\textbf{Detailed Description Task}\quad
The description task aims to evaluate the model's information extraction capabilities, serving as a fundamental aspect of video understanding. Across all three subsets, this task requires providing a free-text answer that describes the selected moment. Furthermore, this task allows for analysis of how the model extracts information and generates answers for subsequent tasks. By examining the model's performance in this task, we gain insights into its ability to extract relevant information and generate meaningful responses.\\
\noindent\textbf{Counter-intuitiveness Reasoning Task}\quad
The reasoning task is designed to test the model's ability to reason about the video, and in the three subsets, this question is Why Humorous, Why creative, and Why counter-intuitive and the answer is a free-text explanation. This task is very difficult and involves the model's deep reasoning ability; it requires the model to give a complete explanation using information from the entire video and its own common sense.\\
\noindent\textbf{Higher Level Tasks}\quad
In addition to the three main tasks, we design higher-level tasks to enhance the model's inference abilities on counter-intuitive videos. \textit{Title Task} in HumorQA and CreativeQA requires generating a concise title summarizing the video's content. \textit{Creative Scoring Task} in CreativeQA involves rating the creativity of videos between 0 and 20. \textit{Magic Method Task} in MagicQA requires the model to explain clearly the rationale behind the magic, and its purpose is to test the model's ability to reason more deeply. To ensure the accuracy of the answers, this task is only partially annotated and appears only in the test set, details of which can be found in Appendix \ref{appendix_a:dataset con}.
\\
In addition, some visual QA benchmarks \cite{Whoop} have begun to adopt free-text answers; this flexible answer format enables multimodal LLMs to go beyond multiple choice options (Large Language Models)~\cite{2023videochat,maaz2023videochatgpt}; for Video QA, however, free-text %while multiple-choice and open-ended formats continue to dominate the Video QA dataset.
are less common. To bridge this gap, we also consider some \textbf{free-text answer} evaluations, demanding enhanced video comprehension capabilities from the model. It will also assess the model's ability to generate reasonably lengthy textual responses.
% \input{table/construction}



\subsection{Dataset Statistics}
% FunQA contains \textbf{4,365} counter-intuitive video clips and \textbf{311,950} question-answer pairs, the total length of these videos is \textbf{23.9} h and the average length of video clips is \textbf{19} seconds. The distribution of the video lengths can be seen in Figure \ref{fig:fun_len}. These videos were divided into train / val / test set and the exact quantities can be seen in Figure \ref{fig:fun_split}. FunQA consists of three fine-grained sub-datasets, each sub-dataset contains four well-designed tasks, the distribution of specific tasks and their corresponding QApairs can be seen in Figure \ref{fig:fun_qa}. Figure \ref{} shows the word cloud map of FunQA as a whole, while Figures \ref{fig:wc_h}, \ref{fig:wc_c} and \ref{fig:wc_m} show the word cloud maps of the three sub-data sets, respectively., the other detailed statistics of sub-datasets are given in Table \ref{T:main} and Appendix \ref{section: FunQA statics}.\\ 
% FunQA has a well-established annotation process and high annotation quality, our cross-validation statistics can be seen in Figure \ref{fig:fun_val}, and the average length of our answers is up to -words, the specific data can be seen in Figure \ref{fig:fun_avg}.Besides, the heat map of the distribution of the results of the Counter-intuitive Timestamp Localization Task over the whole video can be seen in Figure \ref{fig:fun_moment}, which shows that the labeled results are in line with common sense.
% % % static table----------------------------
\begin{table}[!t]
\centering
\caption{\textbf{Statistics of the FunQA dataset.}
\label{T:statics}}
\resizebox{0.8\linewidth}{!}{

\begin{tabular}{lccccccc}
\hline
\multirow{2}{*}{Datase} & \multicolumn{4}{c}{Videos} & \multirow{2}{*}{Avg len(s)} & \multirow{2}{*}{Total len(h)} & \multirow{2}{*}{\# QA Pairs} \\ \cline{2-5}
           & Train & Val & Test & Total &    &       &       \\ \hline
HumorQA    & 1,244  & 350 & 172  & 1,769  & 7  & 3.72  & 141,250  \\
CreativeQA & 657   & 185 & 85   & 927   & 48 & 12.53 & 78,740  \\
MagicQA    & 1,171  & 334 & 167  & 1,672  & 16 & 7.65  & 91,960  \\ \hline
FunQA      & 3,072  & 869 & 424  & 4,365  & 19 & 23.90 & 311,950 \\ \hline \\
\end{tabular}
}
\end{table}
%static table ----------------------------------
% \\
% The distribution of question word length is shown in \ref{fig:fun_avg}(a). Questions are on average 34.24 words, which is much longer than existing VideoQA datasets (e.g., 8.7 in Activity-QA \cite{yu2019activitynet}, 11.6 in NExT-QA \cite{xiao2021nextqanext} ). And the median of each subtask is slightly less than the mean. The distribution of Moment is shown in \ref{fig:fun_moment}(b). We can clearly find that the distribution of counter-intuitive emotions in three different videos is very different. Humor, Creative and Magic videos account for 75\%, 23\% and 61\% of the videos respectively. HumorQA, Creative, Magic QA word cloud is shown in \ref{fig:wc_h, fig:wc_h, fig:wc_m}(c), FunQA full word clouds are in \ref{Data Statistics}.
FunQA contains \textbf{4,365} counter-intuitive video clips and \textbf{311,950} question-answer pairs, the total length of these videos is \textbf{23.9h} and the average length of video clips is \textbf{19} seconds. FunQA consists of three fine-grained subsets, each one containing well-designed tasks. The specific numbers and split of videos can be seen in Figure \ref{fig:dataset statistics} (d). Each subset's video lengths can be seen in Figure \ref{fig:dataset statistics} (d). The specific number of QA pairs for each task can be seen in Figure \ref{fig:dataset statistics} (e). \\
For our localization task, the timestamp heat map for the three different types of videos can be seen in Figure~\ref{fig:dataset statistics} (g), which shows the high-frequency time span of the answer. For the description and reasoning tasks, the average length of the words in their free-text answers reached \textbf{34.24}, which is much longer than existing VideoQA datasets (e.g., 8.7 in Activity-QA~\cite{yu2019activitynet}, 11.6 in NExT-QA~\cite{xiao2021nextqanext}). The specific word count of each task is shown in Figure \ref{fig:dataset statistics} (h). FunQA has a well-established annotation process and high annotation quality, the result of our annotation consensus evaluation can be seen in Figure \ref{fig:dataset statistics} (i). For each video category, more than 90\% of the annotations exhibit a high level of consensus, with only 1\% of the content showing low consensus. Approximately 8\% of the data demonstrates variations in consensus, thus highlighting the objectivity of our dataset. \\
HumorQA, CreativeQA, and MagicQA word clouds are shown in Figure \ref{fig:dataset statistics} (a-c). More statistics and FunQA full word cloud are given in Appendix \ref{appendix_a: Dataset statics}.

  
\subsection{Dataset Construction Pipeline} 
\label{Section: dataset construction pipeline}
FunQA dataset construction pipeline was in three stages: Pre-processing, Manual Annotation, and Post-Processing. The whole process took about 900 hours with over 50 highly educated undergraduates as part-time annotators, and we paid crowd-workers a target of ¥25/hr. More details of the dataset construction pipeline can be seen in Appendix \ref{appendix_a:dataset con}.\\
\noindent\textbf{Pre-processing}\quad
Initially, we crawled videos from TikTok, Bilibili, and YouTube (these videos are a collection of surprising videos). Then we performed a two-stage manual cleaning and cutting process on the collected videos to ensure counter-intuitive features and video quality and to exclude non-ethical content and sensitive information, resulting in a cut video clip.\\
\noindent\textbf{Manual Annotation}\quad
We annotated the videos according to the characteristics of different tasks design in Chinese. We screen and train the annotators to ensure the accuracy and high quality of the annotation, and finally produce the original annotated files. After the first round of annotation, we conducted a secondary annotation of 10\% of the tasks and performed Consensus Evaluation to ensure the objectivity of our annotations.\\
\noindent\textbf{Post-processing}\quad 
Based on our carefully designed tasks and high-quality annotations, we expanded our dataset using GPT-3.5. Firstly, we automatically translated the Chinese annotations into English. Subsequently, we generated more QA pairs that were faithful to the original ideas but presented differently. This not only made FunQA multilingual but also increased the number of QA pairs to 312K. Additionally, we created diverse task types, such as FunQA multiple-choice and FunQA dialogue. More details are given in Appendix \ref{appendix_a: Augmentation} and Appendix \ref{appendix_b: FunQA Extended dataset}.


% 

% Please add the following required packages to your document preamble:
% \usepackage{booktabs}
% \usepackage{multirow}
% \usepackage[table,xcdraw]{xcolor}
% If you use beamer only pass "xcolor=table" option, i.e. \documentclass[xcolor=table]{beamer}
\begin{table}[!t]
\caption{\textbf{Comparison between FunQA and other existing benchmarks.}
Compare to other datasets, FunQA revolves around the captivating realm of interesting and counter-intuitive videos. The tasks within FunQA are specifically designed to challenge the vision capabilities of models, requiring strong skills in producing an in-depth description, interpretation, and spatial-temporal reasoning.
% FunQA's domain is interesting and counter-intuitive videos, FunQA selects videos that are humanly and repeatedly screened for high quality, FunQA's tasks focus on vision, and FunQA's answers require strong description, interpretation, and spatial-temporal reasoning skills already powerful in Free-text generation. FunQA is therefore designed to look at deeper description and explanation. 
Here we clarify the abbreviation in the table. \textbf{Anno.}: Annotation; \textbf{M}: Manual, \textbf{A}: Automatic; For Input, \textbf{V, A, S,} and \textbf{B} denote Video, Audio, Subtitle, and Bounding-box; \textbf{VC} means visual-centric, \textbf{Desc.} means Description, \textbf{Expl} for Explanation, \textbf{STR} for Spatial-temporal Reasoning. For QA Tasks, \textbf{MC} denotes Multiple Choice QA, \textbf{OE} means Open Ended QA, and \textbf{FT} means Free Text QA.
\label{T:main}}
\vspace{-0.05mm}
\resizebox{1.0\linewidth}{!}{
\begin{tabular}{@{}
>{\columncolor[HTML]{FFFFFF}}l 
>{\columncolor[HTML]{FFFFFF}}l 
>{\columncolor[HTML]{FFFFFF}}c 
>{\columncolor[HTML]{FFFFFF}}c 
>{\columncolor[HTML]{FFFFFF}}c 
>{\columncolor[HTML]{FFFFFF}}c 
>{\columncolor[HTML]{FFFFFF}}c 
>{\columncolor[HTML]{FFFFFF}}c 
>{\columncolor[HTML]{FFFFFF}}c 
>{\columncolor[HTML]{FFFFFF}}c 
>{\columncolor[HTML]{FFFFFF}}c 
>{\columncolor[HTML]{FFFFFF}}c @{}}



\toprule
\cellcolor[HTML]{FFFFFF}     
& \cellcolor[HTML]{FFFFFF}                           
& \cellcolor[HTML]{FFFFFF} 
& \multicolumn{3}{c}{\cellcolor[HTML]{FFFFFF}Video} & \multicolumn{5}{c}{\cellcolor[HTML]{FFFFFF}Question Answer}                         \\ 

\cmidrule(l){4-6} \cmidrule(l){7-12} 

\multirow{-2}{*}{\cellcolor[HTML]{FFFFFF}Datasets} 
& \multirow{-2}{*}{\cellcolor[HTML]{FFFFFF}Domain}                                        
& \multirow{-2}{*}{\cellcolor[HTML]{FFFFFF}Anno.} 
& Avg length (s) & \# Clips (K) & Input                                                           
& \# QA pairs (K) & VC & Desc. & Expl. & STR & QA Task      

\\ \cmidrule(r){1-12}

MarioQA \cite{mun2017marioqa} 
& Games  & M & 3-6 & 188 & V, A & 188 & Yes & No & Yes & Yes & OE \\
TGIF-QA \cite{jang-IJCV-2019} 
& Social Media & M & 3.1 & 71.7 & V & 165.2 & Yes & Yes & No & Yes & MC \& OE \\
MovieQA \cite{tapaswi2016movieqa}                                            
& Movies & A & 202.7 & 6.77 & V, A, S & 6.4 & No & No & Yes & Yes & MC \\
CLEVRER \cite{yi2020clevrer}                                           
& Synthetic Video & M & 5 & 20 & V & 305 & Yes & No & Yes & Yes & OE \\
TVQA \cite{lei2018tvqa}    
& TV shows & M & 76.2 & 21.8 &V, A & 152.5 & No & No & Yes & No & MC \\
TVQA+ \cite{lei2019tvqa}  
& TV shows & M &7.2 & 4.2 & V, A, B & 29.4 & No & No & Yes & Yes & OE \\
Social-IQ \cite{zadeh2019social}                                          
& Web videos & M & 99 & 1.25 & V, A, S, B & 7.5 & Yes & No & Yes & No & MC \\
NExT-QA \cite{xiao2021nextqanext}                                            
& Daily life & M & 44 & 5.4 & V, A & 52 & Yes & Yes & Yes & Yes & MC \& OE \\
KnowIT-VQA \cite{garcia2020knowit}                                        
& TV shows & M & 60 & 12 & V, A, S & 24 & Yes & No & Yes & Yes & MC \\
AGQA \cite{grundemclaughlin2021agqa}                                               
& Social Media & A & 30 & 9.6 & V, A & 192 & Yes & No & Yes & Yes & OE \\
AVQA \cite{yang2022avqa}                                               
& Social Media & M & 60 & 9.3 & V, A & 57.3 & Yes & No & Yes & Yes & MC \\
STAR \cite{wu2021star}                                               
& Daily life & A & - & 22 & V & 60 & Yes & No & No & Yes & MC  \\
Env-QA \cite{Gao_2021_ICCV}                                             
& Egocentric Video & M & 20 & 23.3 & V & 85.1 & Yes & No & No & Yes & MC \\
FIBER \cite{castro2022fiber}                                              
& Daily life & M & 10 & 28 & V, A, S & 2 & No & Yes & No & Yes & OE \\
\midrule

HumorQA                                           
& Daily life & M & 7 & 1.8 & V, A & 141.3 & Yes & Yes & Yes & Yes & OE \& FT \\

CreativeQA                                     
& Performance & M & 48 & 0.9 & V, A & 78.7 & Yes & Yes & Yes & Yes & OE \& FT \\

MagicQA                                  
& Magic & M & 10 & 1.6 & V, A & 91.9 & Yes & Yes & Yes & Yes & OE \& FT \\
\midrule
\textbf{FunQA} 
& \textbf{Surprising videos} 
& M & 19 & 4.3 & V, A & \textbf{312} & Yes & Yes & Yes & Yes & \textbf{OE \& FT} \\ \bottomrule 
\end{tabular}}
\vspace{-0.5cm}
\end{table}


% \subsection{Dataset Comparision}

% As shown by the benchmarks in Table \ref{T:main}, there is growing interest in using question responses to measure video reasoning ability \cite{tapaswi2016movieqa, lei2018tvqa, lei2019tvqa, mun2017marioqa, jang-IJCV-2019, yi2020clevrer, xiao2021nextqanext, garcia2020knowit, grundemclaughlin2021agqa, wu2021star, Gao_2021_ICCV, castro2022fiber}. However, the data sources for these well-known benchmarks are often generated semi-automatically from well-known TV shows or films \cite{tapaswi2016movieqa, lei2018tvqa, lei2019tvqa}, so their tasks and answers often rely on dialogue and plot summaries rather than the visual content of the video \cite{tapaswi2016movieqa, lei2018tvqa, lei2019tvqa, kim2017deepstory, zadeh2019social}, leading to models that rely more on dialogue than on visual input, thus reducing the validity of the benchmarks for measuring visual reasoning \cite{tapaswi2016movieqa, lei2018tvqa, lei2019tvqa}.
% An increasing number of benchmarks are now often designed specifically to surprise by focusing on enhancing a particular domain of the model. \cite{yang2022avqa} focuses on improving the model's ability to process audio during problem solving. Moreover, \cite{xiao2021nextqanext, grundemclaughlin2021agqa, wu2021star}  currently focus on improving spatio-temporal reasoning.
% Other pure video benchmarks suffer from biases and simplicity associated with human-generated questions \cite{yu2019activitynetqa, tapaswi2016movieqa, jang-IJCV-2019, lei2018tvqa} or descriptions \cite{xu2017video, zeng2017leveraging}. The largest human annotated \cite{lei2018tvqa} and generated \cite{maharaj2017dataset} datasets contain 152.5K and 349K questions. Picture understanding has recently been increasingly studied in the context of strange, bizarre, and anomalous understanding, but there are no relevant datasets or benchmarks in the field of video understanding. FunQA surpasses the current VideoQA datasets in several aspects:\\
%  1) FunQA distinguishes itself through its diverse video types, each presenting unique characteristics within its respective subset. HumorQA involves incongruity and common sense elements, CreativeQA showcases intricate imitation behavior, and MagicQA entails the impossibility. FunQA's videos are longer and more informative, providing ample visual content for analysis.  \\
% 2) Secondly, FunQA places a strong emphasis on inference in its task design. This necessitates the model to possess deeper inference capabilities and leverage spatial and temporal information effectively. \\
% 3) Among the available VideoQA datasets, FunQA stands out for its inclusion of free-text type answers. Notably, the average answer length in Detailed Description Task and Counter-intuitiveness Reasoning Task is significantly longer compared to existing VideoQA datasets. This demands a higher output capability from the model and allows for evaluation using complex metrics. \\
% 4) FunQA benefits from high-quality human annotation, ensuring the reliability and accuracy of the dataset. They allow us to extend the FunQA dataset with GPT to suit different needs, e.g. FunQA Dialog is more in line with the LLM dialogue format, FunQA Multi-choice has proven evaluation metrics and testing is not limited to LLM, any model can be tested on it.


